\section{Morfismos submersivos y universalmente submersivos}
\label{IX.2}

\begin{definition}
\label{IX.2.1}
Un morfismo
\ifthenelse{\boolean{orig}}
{$g\colon S'\to S$}
{$f\colon S'\to S$}
de preesquemas se dice \emph{submersivo}
\index{submersivo (morfismo)|hyperpage}%
si es sobreyectivo y hace de~$S$ un espacio topológico cociente
de~$S'$ (es decir, toda parte~$U$ de~$S$ tal que $f^{-1}(U)$ sea
abierta es abierta). Se dice que~$f$ es \emph{universalmente
submersivo}
\index{universalmente submersivo (morfismo)|hyperpage}%
si, para todo morfismo~$T\to S$, el morfismo
$f'\colon T'=S'\times_S T\to T$ deducido de~$f$ por cambio de base es
submersivo.
\end{definition}

Es inmediato que el compuesto de dos morfismos submersivos
(respectivamente universalmente submersivos) es submersivo
(respectivamente universalmente submersivo), y que un cambio de base de un
morfismo universalmente submersivo da un morfismo universalmente submersivo
(pues la definición está hecha precisamente para ello). Si~$fg$ es
submersivo (respectivamente universalmente submersivo), también lo es~$f$.

\begin{exemples}
\label{IX.2.2}
a) Un morfismo sobreyectivo que sea abierto o cerrado es submersivo; por
tanto, un morfismo sobreyectivo universalmente cerrado o universalmente
abierto es universalmente submersivo. Por ejemplo, \emph{un morfismo propio
sobreyectivo es universalmente submersivo}. Por otra parte, \emph{un
morfismo fielmente plano y cuasicompacto es universalmente submersivo}
(\Ref{VIII.4.3}). Estos serán para nosotros los dos casos más importantes.
\end{exemples}

Pueden aplicarse
\marginpar{232}
a un morfismo submersivo o universalmente submersivo~$g\colon S'\to S$ los
razonamientos de~\Ref{VIII.4.3}; se obtiene, en particular:
\begin{proposition}
\label{IX.2.3}
Supongamos que~$g\colon S'\to S$ es submersivo. Entonces el siguiente
diagrama de aplicaciones es exacto:
$$
\xymatrix@C=.5cm{\Ouv (S)\ar[r] &
\Ouv (S')\ar@/^2pt/[r]\ar@/_2pt/[r] & \Ouv(S''),}
$$
donde~$S''=S'\times_S S'$ y~$\Ouv(X)$ designa el conjunto de las partes
abiertas del preesquema~$X$.
\end{proposition}

\begin{proposition}
\label{IX.2.4}
Sean~$g\colon S'\to S$ un morfismo universalmente submersivo,
$f\colon X\to Y$ un $S$-morfismo y~$f'\colon X'\to Y'$ el
$S'$-morfismo que se deduce de él por cambio de base. Para que~$f$ sea
\ifthenelse{\boolean{orig}}
{abierta (respectivamente cerrada)}
{abierto (respectivamente cerrado)}, basta con que~$f'$ lo sea. Para que~$f$
sea \ifthenelse{\boolean{orig}}
{universalmente abierta, respectivamente universalmente cerrada,
respectivamente separada}
{universalmente abierto, respectivamente universalmente cerrado,
respectivamente separado}, es necesario y suficiente que~$f'$ lo sea. Si,
además, $g$ es cuasicompacto y~$f$ localmente de tipo finito, para que~$f$
sea propio es necesario y suficiente que~$f'$ lo sea.
\end{proposition}

Para este último punto, obsérvese que, si~$f'$ es propio y, por tanto,
cuasicompacto, entonces~$f$ es cuasicompacto (\Ref{VIII.3.3}) y, por
consiguiente, de tipo finito, puesto que es localmente de tipo finito. Por
otra parte, es separado y universalmente cerrado por lo que precede; luego
es propio.

\begin{proposition}
\label{IX.2.5}
Sea~$S'$ un preesquema de tipo finito sobre el espectro~$S$ de un anillo
local noetheriano completo. Supongamos que la fibra del punto cerrado~$s$
de~$S$ sea finita, de modo que los anillos locales de~$S'$ en los puntos
$s'$ de esta fibra sean finitos sobre~$A=\cal{O}_s$. Sea~$S''$ el esquema
suma disjunta de los espectros de los~$\cal{O}_{S',s'}$ en cuestión,
considerado como $S$-esquema finito. Para que~$g\colon S'\to S$ sea
universalmente submersivo es necesario y suficiente que el morfismo
estructural~$S''\to S'$ sea sobreyectivo.
\end{proposition}

Como existe un $S$-morfismo natural~$S''\to S'$ y todo morfismo finito
sobreyectivo es universalmente submersivo por~\Ref{IX.2.2}, la condición
enunciada es suficiente. Demostremos, pues, que, si~$S''\to S'$ no es
sobreyectivo, entonces~$g$ no es universalmente submersivo. En efecto, sea
$t$ un punto de~$S$ que no pertenece a la
\marginpar{233}
imagen de~$S''$; existe entonces
un $S$-esquema~$T$, espectro de un anillo de valoración discreta, cuya
imagen en~$S$ es~$\{s,t\}$. Obsérvese que la imagen de~$S''$ en~$S'$ es
abierta, pues el morfismo~$S''\to S'$ es un isomorfismo local; además, esta
imagen contiene~$S'_s$ y no corta a~$S'_t$. Se sigue que la preimagen de
esta última en~$T'=S'\times_S T$ es \emph{abierta} y coincide con la
preimagen del punto cerrado de~$T$. Esto
demuestra que~$T'\to T$ no es submersivo y, por tanto, que~$S'\to S$ no es
universalmente submersivo.

\begin{remarque}
\label{IX.2.6}
Utilizando el criterio~\Ref{IV.6.3} para que una parte constructible de un
espacio noetheriano sea abierta, se obtiene fácilmente el siguiente criterio
valuativo para que un morfismo~$g\colon S'\to S$ \emph{de tipo finito}, con
$S$ localmente noetheriano, sea universalmente submersivo: es necesario y
suficiente que, para todo $S$-esquema~$T$ que sea el espectro de un anillo de
valoración discreta, poniendo~$T'=S'\times_S T$, la preimagen en~$T'$ del
punto cerrado de~$T$ no sea abierta.
\end{remarque}
